%------------------------------------------------------------------------------%
\documentclass[fleqn,utf8,aspectratio=169,12pt]{beamer}

\usepackage{calculo_varias_variaveis-1}

\title{Diferenciabilidade}

%------------------------------------------------------------------------------%
\begin{document}

\addtitle

%------------------------------------------------------------------------------%
\section{Diferenciabilidade}

%------------------------------------------------------------------------------%
\begin{frame}{Derivada Ordinária -- Uma variável}

  Variação da função com relação à variável

  \pause
  \vspace{\baselineskip}
  Derivada da função \(f\colon R\subset\R\to\R\) no ponto $a$

  \[
    \frac{df}{dx}(a)
    =\; \lim_{h\to 0} \frac{f(a+h)-f(a)}{h}
  \]
\end{frame}

%------------------------------------------------------------------------------%
\begin{frame}{Diferenciabilidade em 1D}
	Se a \emph{derivada de $f$ existe} em $a$ então $f$ é \emph{diferenciável} em $a$

	\pause
	\vspace{0.3\baselineskip}
	Existe uma \emph{reta tangente} a $f$ no ponto \(\big(a, f(a)\big)\)

	\pause
	\vspace{0.3\baselineskip}
	Podemos \emph{aproximar} a função pela reta em uma vizinhança de $a$
  \[
    \pause
    \Delta x = x - a
  \]
  \[
    \pause
    \Delta y = f(x) - f(a)
  \]
  \[
    \pause
    \hphantom{\Delta y} = f'(a)\Delta x + \emph{\varepsilon \Delta x}
  \]
  \[
    \pause
    \emph{\varepsilon\to 0} \;\text{ quando }\; \emph{\Delta x \to 0}
  \]
\end{frame}

%------------------------------------------------------------------------------%
\begin{frame}{Definição -- Derivadas Parciais}
  As derivadas parciais, com relação a \emph{$x$} e \emph{$y$},
  da função \(f\colon R\subset\R^2\to\R\),
  no ponto \((x, y) = (a, b)\)

  \[
    \D{f}{x}(a, b)
    =\; \lim_{h\to 0} \frac{f(\emph{a+h}, b)-f(\emph{a}, b)}{h}
  \]

  \[
    \D{f}{y}(a, b)
    =\; \lim_{h\to 0} \frac{f(a, \emph{b+h})-f(a, \emph{b})}{h}
  \]

  Desde que os limites existam
\end{frame}

%------------------------------------------------------------------------------%
\begin{frame}{Diferenciabilidade em $n$ Dimensões}
	A existência das derivadas parciais \emph{não} garante diferenciabilidade

	\pause
	\vspace{\baselineskip}
	Diferenciabilidade está associada a existência do \emph{plano tangente}
\end{frame}

%------------------------------------------------------------------------------%
\begin{frame}{Definição: Diferenciabilidade em 2D}
  Uma função \(z=f(x, y)\) é \emph{diferenciável} em \((a, b)\)

  \pause
  \vspace{0.5\baselineskip}
  se \(\D{f}{x}(a, b)\) e \(\D{f}{y}(a, b)\) existem

  \pause
  \vspace{0.5\baselineskip}
  e \(\Delta z = f(a+\Delta x, b+\Delta y)-f(a, b)\)

  \pause
  \vspace{0.5\baselineskip}
  satisfaz
  \[
    \Delta z
    \;=\; f_x(a, b) \Delta x
    \;+\; f_y(a, b) \Delta y
    \;+\; \emph{\varepsilon_1 \Delta x}
    \;+\; \emph{\varepsilon_2 \Delta y}
  \]
  \[
    \pause
    \emph{\varepsilon_1, \varepsilon_2 \to 0}
    \;\text{ quando }\;
    \emph{\Delta x, \Delta y \to 0}
  \]
\end{frame}

%------------------------------------------------------------------------------%
\begin{frame}{Teorema do Incremento}
	Suponha que as derivadas parciais de primeira ordem de $f(x, y)$ sejam
	definidas em uma região aberta $R$, que contém o ponto $(a, b)$
  \pause
	e que \emph{$f_x$ e $f_y$ sejam contínuas em $(a, b)$},
	\pause
	então a variação
	\[
		\pause
    \Delta z
		\;=\; f\left(a+\Delta x, b + \Delta y\right)
		\;-\; f(a, b)
	\]
	\pause
  com \(\left(a+\Delta x,\, b + \Delta y\right)\in R\),
	\pause
  satisfaz
	\[
		\Delta z
		\;=\; f_x(a, b) \Delta x
		\;+\; f_y(a, b) \Delta y
		\;+\; \emph{\varepsilon_1 \Delta x}
		\;+\; \emph{\varepsilon_2 \Delta y}
	\]
	\[
		\pause
    \emph{\varepsilon_1, \varepsilon_2 \to 0}
    \;\text{ quando }\;
    \emph{\Delta x, \Delta y \to 0}
	\]
\end{frame}

%------------------------------------------------------------------------------%
\begin{frame}{Corolário do Teorema do Incremento}
  Se \(\D{f}{x}(a, b)\) e \(\D{f}{y}(a, b)\) existem

  \pause
  \vspace{0.5\baselineskip}
  e são \emph{contínuas} em uma região aberta $R$

  \pause
  \vspace{0.5\baselineskip}
  então $f$ é \emph{diferenciável} em $R$
\end{frame}

%------------------------------------------------------------------------------%
\begin{frame}{Diferenciabilidade em 2D}
  Dizemos que $f$ é \emph{diferenciável}
  se ela é diferenciável em todos os
  pontos do seu domínio

  \pause
  \vspace{0.5\baselineskip}
  Nesse caso dizemos que seu gráfico é uma \emph{superfície lisa}
\end{frame}

%------------------------------------------------------------------------------%
\begin{frame}{Teorema Diferenciabilidade e Continuidade}
    Se uma função $f(x, y)$ é \emph{diferenciável} em \((a, b)\)

    \pause
    \vspace{0.5\baselineskip}
    então ela é contínua em \((a, b)\)
\end{frame}

%------------------------------------------------------------------------------%
\section{Lista Mínima}

%------------------------------------------------------------------------------%
\begin{frame}{Lista Mínima}
  Cálculo Vol. 2 do Thomas 12\textsuperscript{a} ed. -- Seção 14.3
  \begin{enumerate}
    \item Estudar o texto da seção
%    \item Resolver os exercícios:
  \end{enumerate}

  \vfill
  Atenção: A prova é baseada no livro, não nas apresentações
\end{frame}

%------------------------------------------------------------------------------%
\end{document}
%------------------------------------------------------------------------------%
